\begin{definition}[Spectral Leap]
Define the spectral gap of $\mathcal L_E$ by
\[
\lambda_1(E) :=
\inf_{\phi \perp \ker \mathcal L_E}
\frac{\langle \mathcal L_E \phi, \phi \rangle}{\|\phi\|^2}.
\]
\end{definition}

\begin{lemma}[Spectral Coercivity]
The condition $\lambda_1(E) > 0$ is equivalent to the existence of a constant
$c(E)>0$ such that
\[
\hat h(P) \ge c(E)\,\|P\|^2
\quad \text{for all } P \in E(\mathbb{Q}) \setminus E(\mathbb{Q})_{\mathrm{tors}}.
\]
\end{lemma}

